\section{Обговорення}
\subsection{Практична цінність інтеграції}
Головна перевага Social --- неперервність між інакше відокремленими аналітичними завданнями. Аналітик може переглянути зафіксовану рамку збору, простежити нормалізований запис до доступного матеріалу джерела, порівняти кандидатний граф копій із точним обчисленням, дослідити умовну нульову модель і зберегти нерозв'язані докази в процесі, орієнтованому на звіт. Цінність фреймворку полягає у зменшенні неоднозначності щодо того, який запис, припущення й перетворення створили результат. Економію часу аналітика чи якість кінцевого розслідування порівняно з роботою лише експертів не вимірювали, тож ці переваги є архітектурними можливостями, а не встановленими ефектами продуктивності.

Кілька проєктних рішень навмисно скромні. JSONL і SQLite придатні для перегляду й практичні для одного дослідника. Детерміновані правила легше відтворювати, ніж необмеженого агента, а локальна обробка медіа дозволяє не надсилати кожен приватний актив до хмарної моделі. Однак прозорість правила не робить його семантично правильним: контрприклад із числовим присудком є прямою ілюстрацією, а звуження цього правила після аудиту не перетворило його на валідовану модель випливання. Найбезпечніший загальний запускач наразі утримується від прийняття; амбітніший кураторський шлях показує, наскільки оманливою може стати структурна повнота, коли міткам прийняття довіряють надто легко.

\subsection{Відповіді на дослідницькі питання}
ДП1: отримання даних, знімки й типізований офлайн-процес інтегровано частково, тоді як обробка медіа, кілька збирачів і кураторські збирачі запускаються окремо. ДП2: інваріанти посилань, формату хешів, граничних часів і заміщення забезпечуються; випливання, автентичність і незалежність джерел --- ні. ДП3: точний граф, збережені підсумки рандомізації та префікси із зерном відтворюються на зафіксованих входах. ДП4: комунікаційні й фізичні часові шкали розділені, але фізичні інтервали надає аналітик. ДП5: збої числового присудка й заперечення проходили програмні бар'єри, доки правило не звузили, а ширша семантична помилка лишається невиміряною. ДП6: обмежений процес рівня~2 підтримується, незалежно валідоване розслідування --- ні.

\subsection{Уроки для доказово обґрунтованого програмного забезпечення}
По-перше, походження є набором тверджень, що потребують перевірки, а не просто колекцією хешів. Хешований документ ліній походження може приписати результат не тому корпусу. По-друге, надійність багатовимірна: успіх тестів, правильна арифметика, цілісність цитувань, семантичне випливання й істинність подій можуть розходитися. По-третє, автономність слід приписувати конкретному шляху виконання та завданню. Відеопілот, ручний перелік подій і знайдений набір уривків не можна риторично об'єднати в незалежно валідованого автономного дослідника. По-четверте, чесні «невідомо» є корисними результатами. Вони зберігають розуміння того, яких додаткових доказів бракує, і не дають правдоподібному доповненню замінити відсутнє спостереження.

Наступним дослідницьким пріоритетом є незалежне семантичне оцінювання на відокремлених інцидентах і сім'ях джерел, а не складніша оркестрація моделей. Зафіксований еталонний набір може підтримувати розробку моделей і регресійне оцінювання без потреби вручну перевіряти кожен майбутній інцидент, але твердження про застосування мають лишатися обмеженими продемонстрованим перенесенням. Базова лінія лише з LLM чи кількома моделями потребувала б узгоджених входів, завдань і бюджетів. Такого порівняння тут не виконували, а згода моделей усе одно не дала б незалежного підтвердження.
